% !TEX root = ../../../main.tex
% Continue front-right. Follow Chapter 3 order.
\section{Sequences}
\label{sec:appendix:formula-sheet-sequences}
$x_n\to L$ means
$\forall\varepsilon>0\ \exists N\ \forall n\ge N\quad |x_n-L|<\varepsilon$.
If $x_n\to L$, $y_n\to M$, limit laws give
\[
  \begin{gathered}
    ax_n+by_n\to aL+bM,\\
    x_ny_n\to LM,\qquad x_n/y_n\to L/M\ (M\ne0).
  \end{gathered}
\]
Squeeze: eventual $u_n\le x_n\le v_n$ and $u_n,v_n\to L$ imply $x_n\to L$.

\ifdefined\standaloneformulasheet\else
\formulagroup{Binomial identity}
\[
  \begin{gathered}
    (u+v)^n=\sheetsum_{j=0}^n\binom nj u^{n-j}v^j,\\
    \binom nj=\frac{n!}{j!(n-j)!}.
  \end{gathered}
\]
\fi

\formulagroup{Standard limits}
Bernoulli: $(1+t)^n\ge1+nt$ for $t\ge-1$.
For fixed $c>0$, $a>1$, and $p>0$,
\[
  \begin{gathered}
    n^{-p}\to0,\quad c^{1/n}\to1,\quad n^{1/n}\to1,\\
    n^{p/n}\to1,\quad n^p/a^n\to0,\\
    (a^n+n)^{1/n}\to\max\{1,a\}\quad(a\ge0),\\
    (1+1/n)^n\to e,\qquad (1+1/n)^{1/n}\to1.
  \end{gathered}
\]
$r^n\to0$ for $|r|<1$; $r^n\to+\infty$ for $r>1$.
For a polynomial with positive leading coefficient, $P(n)^{1/n}\to1$.

\formulagroup{Monotonicity, AM--GM, and subsequences}
A bounded increasing sequence tends to $\sup\{x_n\}$ and has infimum $x_1$;
a bounded decreasing one tends to $\inf\{x_n\}$ and has supremum $x_1$.
\[
  (a_1\cdots a_k)^{1/k}\le\frac{a_1+\cdots+a_k}{k}
  \quad(a_i\ge0,\ k\in\N);
\]
equality holds exactly when all $a_i$ agree.
Bounded real sequences have convergent subsequences.
Two distinct subsequence limits rule out convergence.
For indices $n\equiv0,1,2\pmod3$, respectively,
$\sin(2\pi n/3)=0,\sqrt3/2,-\sqrt3/2$.

\formulagroup{Tail bounds and Cauchy}
\[
  \begin{aligned}
    \liminf x_n&=\lim_n\inf_{k\ge n}x_k,\qquad
    \limsup x_n=\lim_n\sup_{k\ge n}x_k,\\
    x_n\to L&\iff\liminf x_n=\limsup x_n=L.
  \end{aligned}
\]
For bounded sequences these are the extreme subsequence limits.
In $\R$, convergence is equivalent to the Cauchy condition
\[
  \forall\varepsilon>0\ \exists N\ \forall m,n\ge N\quad
  |x_m-x_n|<\varepsilon.
\]

